% ============================================================================
% Fichier  : partie6/P06_C26_retex.tex
% Titre    : Retour d'Expérience, Résilience et Mesure
% Auteur   : MINKA MI NGUIDJOÏ Thierry Emmanuel
% Version  : 1.0 -- novembre 2026
% ============================================================================
% Sources : Reason (1997) ; ISO 22361:2022 ; ISO 22301:2019 (10) ; instrument
% OUI / NON / NEK (prépublications HAL de l'auteur). Questionnaire de
% résilience de LiZbiZ : cas fictif ; mu, U, H et intervalles de Wilson à
% 95 % calculés et vérifiés.
% ============================================================================

\chapter{Retour d'Expérience, Résilience et Mesure}
\label{chap:retex}
\chaptermark{Retour d'expérience et mesure}

\epigraph{%
    « Une organisation apprend quand ses erreurs\\
    changent ses pratiques, pas quand elle les raconte. »%
}{Principe du cours}

\niveauChapitre{Expert}{Ensemble du cours.}

\begin{boxobjectifs}
\begin{itemize}
    \item Conduire un retour d'expérience~: faits, écarts, causes, actions.
    \item Remonter aux causes profondes sans désigner de coupable.
    \item Suivre les actions jusqu'à leur effet.
    \item Mesurer la résilience avec l'instrument \OUI{}~/ \NON{}~/ \NEK{}
          et les mesures $\mu$, $U$, $H$.
    \item Dire la précision d'une mesure et ce qu'elle ne permet pas de
          conclure.
\end{itemize}
\end{boxobjectifs}

\bigskip

% ============================================================================
\section*{Question directrice : LiZbiZ est-elle plus résiliente en octobre qu'en mars, et le sait-elle~?}
% ============================================================================

Le 2~avril, la réunion de retour d'expérience de l'Opération GAROUA a
produit quarante idées et aucune décision. En mai, Mme~MOUAFO a repris la
méthode. En octobre, le conseil demande~: « sommes-nous plus résilients~?
De combien~? Avec quelle certitude~? » Ce dernier chapitre répond, et
résume le cours.

\medskip
\textbf{Plan du chapitre.}
\begin{enumerate}
    \item Le retour d'expérience (section~\ref{sec:c26-retex}).
    \item Remonter aux causes (section~\ref{sec:c26-causes}).
    \item Suivre jusqu'à l'effet (section~\ref{sec:c26-suivre}).
    \item Mesurer la résilience (section~\ref{sec:c26-mesurer}).
    \item La précision de la mesure (section~\ref{sec:c26-precision}).
    \item Synthèse du cours (section~\ref{sec:c26-synthese}).
\end{enumerate}

\separateur

% ============================================================================
\section{Le Retour d'Expérience}
\label{sec:c26-retex}
% ============================================================================

\begin{boxdef}
\termecle{Retour d'expérience} (RETEX)~: démarche structurée qui, après un
incident, une crise ou un exercice, établit les faits, les compare à ce
qui était prévu, en recherche les causes et décide des actions qui
changent durablement les pratiques.
\end{boxdef}

\begin{boxprocedure}
\textbf{Conduire un retour d'expérience}
\begin{enumerate}
    \item \textbf{Établir la chronologie} à partir de la main courante et
          des journaux, avant toute interprétation.
    \item \textbf{Comparer} à ce qui était prévu (plans, objectifs,
          procédures)~: les écarts.
    \item \textbf{Rechercher les causes} de chaque écart, jusqu'aux causes
          d'organisation.
    \item \textbf{Décider} d'un petit nombre d'actions, chacune avec un
          responsable, une échéance et un indicateur d'effet.
    \item \textbf{Diffuser} les enseignements à ceux qui doivent changer
          leurs pratiques.
    \item \textbf{Vérifier}, à l'exercice suivant, que les actions ont
          produit leur effet.
\end{enumerate}
\end{boxprocedure}

\begin{boxalert}
\textbf{Violation typique.} Faire du retour d'expérience la recherche d'un
coupable.

\smallskip
\textbf{Règle~:} dans une culture juste (Chapitre~\ref{chap:perception}),
on cherche pourquoi le système a permis l'erreur, et non qui l'a commise.
Le technicien qui a redémarré les serveurs le 5~mars n'avait ni consigne,
ni grille, ni formation~: c'est l'organisation qu'il faut corriger.
\end{boxalert}

% ============================================================================
\section{Remonter aux Causes}
\label{sec:c26-causes}
% ============================================================================

James Reason distingue les \emph{défaillances actives}, commises au
contact du système (un clic, un redémarrage), et les \emph{conditions
latentes}, créées par des décisions d'organisation, parfois anciennes, qui
laissent passer les défaillances~\cite{reason1997}. Le retour
d'expérience remonte des unes aux autres.

\begin{tableaucours}{L{4.4cm}L{7.8cm}}{%
    \tableauHeader{Écart de mars} &
    \tableauHeader{Chaîne des causes}
}
Serveurs redémarrés à 22~h~05 & Pas de consigne~$\rightarrow$ pas de
grille d'escalade~$\rightarrow$ pas de procédure d'incident~$\rightarrow$
aucun dispositif de gestion des risques \\
Cellule réunie à 10~h~30 & Pas de plan de crise~$\rightarrow$ personne ne
savait qui décidait~$\rightarrow$ crise jamais envisagée \\
Onze jours pour facturer & Pas de copie immuable~$\rightarrow$ pas d'objectif
de reprise~$\rightarrow$ pas d'analyse d'impact \\
Compte du prestataire sans second facteur & Exception jamais revue
$\rightarrow$ recommandation de 2023 non suivie~$\rightarrow$ pas de suivi
des audits \\
\end{tableaucours}
\vspace{-0.8em}
\begin{center}{\small\itshape Tableau~26.1. Des défaillances actives aux conditions latentes}\end{center}

% ============================================================================
\section{Suivre jusqu'à l'Effet}
\label{sec:c26-suivre}
% ============================================================================

\begin{boxscenario}{Opération GAROUA. Le retour d'expérience repris en mai}
\textbf{Décisions} (cinq actions au lieu de quarante idées)~: dispositif
de gestion des risques (Partie~II)~; programme d'audit et suivi des
recommandations (Partie~III)~; analyse d'impact et plans (Partie~IV)~;
cellule de crise et fiche réflexe (Partie~V)~; exercices réguliers.

\medskip
\textbf{Effets vérifiés en octobre.} Registre et comité en place~; audit
d'octobre réalisé et suivi~; test de reprise de la facturation en
19~heures~; cellule réunie en 25~minutes à l'exercice. \textbf{Analyse}
(confiance moyenne). Les cinq actions ont produit un effet mesurable~;
leur solidité face à une crise inopinée reste à démontrer.
\end{boxscenario}

% ============================================================================
\section{Mesurer la Résilience}
\label{sec:c26-mesurer}
% ============================================================================

Mme~MOUAFO pose les mêmes vingt questions en mars (après coup, sur
l'état antérieur) et en octobre~: « Existe-t-il un plan de crise
approuvé~? », « Le RTO de la facturation a-t-il été mesuré en test~? »,
« L'annuaire de crise a-t-il été vérifié ce trimestre~? »~\ldots~Chaque
question reçoit \OUI{}, \NON{} ou \NEK{} (connaissance insuffisante pour
conclure), selon l'instrument proposé par l'auteur~\cite{minka2026nek}.

\begin{boxdef}
\textbf{Trois mesures.} Sur $n$ questions, avec $n_O$ réponses \OUI{},
$n_N$ \NON{} et $n_K$ \NEK{}~: $\mu = n_O / n$, la \termecle{conformité}~;
$U = n_K / n$, l'\termecle{incertitude}~; $H = -\sum p \log_2 p$ sur les
trois proportions, l'\termecle{entropie}, au plus $\log_2 3 \approx 1{,}585$
bit.
\end{boxdef}

\begin{boxlecture}{La résilience de LiZbiZ, mars et octobre}
Vingt questions sur la préparation (plans, cellule, exercices, annuaire,
objectifs mesurés, assurance, contrats).

\medskip
\begin{tableaucours}{L{2.6cm}cccc}{%
    \tableauHeader{Date} &
    \tableauHeader{O / N / K} &
    \tableauHeader{$\mu$} &
    \tableauHeader{$U$} &
    \tableauHeader{$H$ (bit)}
}
Mars & 3 / 9 / 8 & 0,15 & 0,40 & 1,458 \\
Octobre & 13 / 3 / 4 & 0,65 & 0,20 & 1,279 \\
\end{tableaucours}

\medskip
\textbf{Constatation.} La conformité passe de 0,15 à 0,65, l'incertitude
de 0,40 à 0,20. \textbf{Analyse} (confiance moyenne). La préparation a
progressé, et l'on sait mieux où l'on en est~: quatre questions restent
sans réponse établie. Compter les \NEK{} comme des \OUI{} afficherait
85~\% en octobre, une image flatteuse que rien ne justifie.

\hfill\lecture{0{,}65}{0{,}20}{1{,}28}
\end{boxlecture}

% ============================================================================
\section{La Précision de la Mesure}
\label{sec:c26-precision}
% ============================================================================

Vingt questions, c'est peu. L'intervalle de Wilson à 95~\% donne, pour
$\mu$, de 0,05 à 0,36 en mars et de 0,43 à 0,82 en octobre~: les deux
intervalles ne se recouvrent pas, la progression de la conformité est
établie. Pour $U$, il va de 0,22 à 0,61 en mars et de 0,08 à 0,42 en
octobre~: les intervalles se recouvrent, la baisse de l'incertitude est
probable mais pas établie avec vingt questions.

\begin{boxalert}
\textbf{Violation typique.} Annoncer au conseil « 65~\% de résilience ».

\smallskip
\textbf{Règle~:} on donne la mesure, son effectif, sa marge et ce qu'elle
mesure (la préparation déclarée et vérifiée, pas la résilience face à une
crise réelle). « Résilience » se démontre par des exercices inopinés et
par les crises elles-mêmes.
\end{boxalert}

\begin{boxmetier}{Gestionnaire de risques}
\textbf{Faire baisser $U$ d'abord.} Mme~MOUAFO consacre le trimestre
suivant aux quatre questions \NEK{}~: chacune deviendra \OUI{} ou \NON{}.
Un \NON{} établi vaut mieux qu'un \NEK{}~: il se traite.
\end{boxmetier}

% ============================================================================
\section{Synthèse du Cours}
\label{sec:c26-synthese}
% ============================================================================

\begin{tableaucours}{L{3.4cm}L{8.8cm}}{%
    \tableauHeader{Partie} &
    \tableauHeader{Ce qu'elle apprend}
}
I. Fondements & Le risque est une incertitude sur des pertes~; nos
jugements sont biaisés~; il faut des rôles, un cadre, et la loi l'exige \\
II. Analyser & Formuler, coter avec prudence, conduire EBIOS RM, chiffrer
avec FAIR, simuler, traiter et rendre compte \\
III. Auditer & Vérifier par la preuve que ce que l'on croit en place l'est,
et suivre jusqu'à l'effet \\
IV. Continuer & Savoir ce qui ne peut pas attendre, préparer les plans,
reconstruire sain, exercer \\
V. Décider en crise & Une cellule organisée, une main courante, des
décisions structurées, une parole tôt et juste \\
VI. Apprendre & Remonter aux conditions latentes, mesurer, et dire ce que
l'on ne sait pas \\
\end{tableaucours}
\vspace{-0.8em}
\begin{center}{\small\itshape Tableau~26.2. Les six parties du cours}\end{center}

\separateur

\begin{boxresume}
\textbf{Ce qu'il faut retenir de ce chapitre~:}
\begin{itemize}
    \item \textbf{Le retour d'expérience}~: chronologie, écarts, causes,
          quelques actions, diffusion, vérification.
    \item \textbf{On remonte} des défaillances actives aux conditions
          latentes, sans chercher de coupable.
    \item \textbf{Une action} n'est finie que lorsque son effet est vérifié.
    \item \textbf{$\mu$, $U$, $H$} mesurent la préparation et ce que l'on en
          ignore~; on fait d'abord baisser $U$.
    \item \textbf{Une mesure} se donne avec son effectif et sa marge.
\end{itemize}
\end{boxresume}

\begin{boxinfo}
\textbf{Points de contrôle de l'élève.} Je sais~: conduire un retour
d'expérience~; distinguer défaillance active et condition latente~;
calculer $\mu$, $U$ et $H$~; lire un intervalle de Wilson~; présenter une
mesure avec sa marge.
\end{boxinfo}

\bigskip

\begin{boxexo}[title={\ding{45}\quad Exercice~26.1. Actif ou latent \niveaubase}]
Classez en défaillance active ou condition latente~: le clic de
l'assistante hameçonnée~; l'absence de politique de dérogation~; le
redémarrage des serveurs.
\end{boxexo}

\begin{boxexo}[title={\ding{45}\quad Exercice~26.2. Calculer \niveaubase}]
Calculez $\mu$, $U$ et $H$ pour 10 \OUI{}, 6 \NON{} et 4 \NEK{}.
\end{boxexo}

\begin{boxexo}[title={\ding{45}\quad Exercice~26.3. Quarante idées \niveaubase}]
Pourquoi la réunion du 2~avril, avec quarante idées, n'a-t-elle rien
changé~? Que fallait-il faire~?
\end{boxexo}

\begin{boxexo}[title={\ding{45}\quad Exercice~26.4. La chaîne \niveaumoyen}]
Construisez la chaîne des causes de l'écart « premier message externe
25~h~30 après la découverte ».
\end{boxexo}

\begin{boxexo}[title={\ding{45}\quad Exercice~26.5. NEK compté \niveaumoyen}]
Calculez la conformité d'octobre si l'on compte les \NEK{} comme des
\OUI{}, puis comme des \NON{}. Pourquoi les deux sont-elles trompeuses~?
\end{boxexo}

\begin{boxexo}[title={\ding{45}\quad Exercice~26.6. Les intervalles \niveaumoyen}]
Expliquez à M.~ETEME, sans formule, pourquoi la progression de la
conformité est établie et celle de l'incertitude seulement probable.
\end{boxexo}

\begin{boxexo}[title={\ding{45}\quad Exercice~26.7. Un indicateur d'effet \niveaumoyen}]
Pour chacune des cinq actions décidées en mai, proposez un indicateur qui
mesure son effet, et non sa seule réalisation.
\end{boxexo}

\begin{boxexo}[title={\ding{45}\quad Exercice~26.8. Faire baisser U \niveaumoyen}]
Pour les quatre questions \NEK{} d'octobre (à imaginer), indiquez la preuve
qui permettrait de répondre \OUI{} ou \NON{}.
\end{boxexo}

\begin{boxexo}[title={\ding{45}\quad Exercice~26.9. Le rapport au conseil \niveauexpert}]
Rédigez la page de Mme~MOUAFO au conseil sur la résilience de LiZbiZ~:
mesures, marges, limites, décisions demandées.
\end{boxexo}

\begin{boxexo}[title={\ding{45}\quad Exercice~26.10. Synthèse \niveauexpert}]
En deux pages, racontez la crise de mars telle qu'elle se serait déroulée
si LiZbiZ avait appliqué ce cours un an plus tôt~: risques, audit,
continuité, crise, retour d'expérience.
\end{boxexo}

\bigskip

\begin{boxinfo}
\textbf{Pour aller plus loin.}
\begin{itemize}
    \item J.~Reason~\cite{reason1997}.
    \item L'instrument \OUI{}~/ \NON{}~/ \NEK{}~\cite{minka2026nek}~;
          l'auditologie~\cite{minka2026auditologie}.
    \item Les UE INF3226 (audit des systèmes d'information), SEC3266
          (sécurité informatique), SEC4031 et SEC4052 (investigation
          numérique), SEC4062 (hacking éthique).
\end{itemize}
\end{boxinfo}

\begin{boxinfo}
\textbf{Conclusion.} Gérer les risques, c'est décider avant de savoir~;
gérer une crise, c'est décider sans avoir le temps de savoir~; apprendre,
c'est savoir enfin, et le dire avec sa marge.
\end{boxinfo}
